\chapter*{Vorwort: Was dieser Band will}
\addcontentsline{toc}{chapter}{Vorwort}
\markboth{Vorwort}{Vorwort}

In Band 15 dieser Reihe, \kursiv{Die Kunst des Zählens}, steht die Graphentheorie als Teil~V: drei Kapitel, vierunddreißig Seiten. Das war zu wenig, und zwar nicht aus Platzgründen, sondern aus einem sachlichen: Die Graphentheorie ist kein Anhängsel der Kombinatorik. Sie ist ein eigenes Gebiet mit eigenen Fragen, eigenen Beweismethoden und einer eigenen Geschichte, die 1736 in Königsberg beginnt und heute in der Wahrscheinlichkeitstheorie, der Algebra und der theoretischen Informatik weitergeht.

Dieser Band holt das nach. Er setzt nichts voraus als die Bereitschaft, Definitionen ernst zu nehmen.

\begin{merksatz}
\boxtitel{Was Sie mitbringen müssen}
\begin{center}
\renewcommand{\arraystretch}{1.5}
\begin{tabular}{@{}P{3.2cm}P{5.8cm}@{}}
\toprule
\textbf{Mengen} & Was $\in$, $\subseteq$, $\cup$, $\cap$ bedeuten. Mehr nicht.\\
\textbf{Summenzeichen} & $\sum$ als abkürzende Schreibweise für \enquote{addiere alles auf}.\\
\textbf{Induktion} & Wird in Kapitel~4 noch einmal vollständig vorgeführt, bevor sie gebraucht wird.\\
\textbf{Matrizen} & Erst ab Kapitel~3, und dort wird erklärt, was eine Matrix ist.\\
\bottomrule
\end{tabular}
\end{center}

\begin{center}
\begin{minipage}{0.88\textwidth}
\kursiv{Wer Band 0 (\enquote{Die Kunst des Rechnens}) durchgearbeitet hat, hat alles. Wer das nicht hat, kommt trotzdem durch --- die vier Dinge oben werden im Text nachgereicht, wo sie zum ersten Mal auftreten.}
\end{minipage}
\end{center}
\end{merksatz}

\section*{Warum die Graphentheorie ein guter Einstieg ist}

Es gibt einen Grund, warum Vorlesungen über Graphentheorie oft im ersten Semester liegen, und er ist nicht, dass der Stoff leicht wäre. Er ist dieser: Die Gegenstände sind endlich, anschaulich und trotzdem tief.

Ein Graph ist nichts als eine endliche Menge von Punkten und eine Auswahl von Paaren dieser Punkte. Man kann ihn auf ein Blatt Papier zeichnen. Man kann ihn abzählen. Es gibt keine Grenzwerte, keine unendlichen Prozesse, keine Konvergenzfragen. Und dennoch enthält dieses schlichte Gebilde Probleme, an denen die besten Mathematiker Jahrzehnte gescheitert sind --- der Vierfarbensatz brauchte hundertvierundzwanzig Jahre, und die Hadwigersche Vermutung von 1943 ist bis heute offen.

\begin{merksatz}
\boxtitel{Was das für Sie bedeutet}
\begin{center}
\begin{minipage}{0.9\textwidth}
Sie können in diesem Gebiet \kursiv{jeden einzelnen Beweis vollständig nachvollziehen}, ohne irgendwo auf Vertrauen angewiesen zu sein. Und Sie stehen nach fünfzehn Kapiteln vor offenen Fragen, die Sie verstehen --- nicht vor offenen Fragen, von denen Sie gehört haben.
\end{minipage}
\end{center}
\end{merksatz}

\section*{Drei Warnungen}

\begin{warnung}
\boxtitelr{Erstens: Das Bild ist nicht der Graph}
Graphen werden gezeichnet, und die Zeichnung hilft. Sie führt aber auch in die Irre, weil derselbe Graph auf unendlich viele Weisen gezeichnet werden kann --- als übersichtliches Fünfeck oder als Knäuel, in dem sich Linien kreuzen.

\begin{center}
\begin{tikzpicture}[scale=0.95]
\begin{scope}
\foreach \i in {0,...,4} {
  \node[knoten] (a\i) at (90+72*\i:0.95) {};
}
\foreach \i/\j in {0/1,1/2,2/3,3/4,4/0} \draw[kante] (a\i) -- (a\j);
\node[font=\scriptsize\itshape, color=grau] at (0,-1.5) {so gezeichnet};
\end{scope}
\begin{scope}[xshift=4.6cm]
\node[knoten] (b0) at (0,0.9) {};
\node[knoten] (b1) at (1.1,0.1) {};
\node[knoten] (b2) at (-0.2,-0.7) {};
\node[knoten] (b3) at (0.9,-0.9) {};
\node[knoten] (b4) at (-1.1,0.2) {};
\draw[kante] (b0) -- (b1);
\draw[kante] (b1) -- (b2);
\draw[kante] (b2) -- (b4);
\draw[kante] (b4) -- (b3);
\draw[kante] (b3) -- (b0);
\node[font=\scriptsize\itshape, color=grau] at (0,-1.5) {oder so};
\end{scope}
\end{tikzpicture}
\end{center}

Beides ist derselbe Graph: fünf Ecken, jede mit genau zwei Nachbarn, alles in einem Kreis. Rechts kreuzen sich zwei Linien --- das ist eine Eigenschaft der \kursiv{Zeichnung}, nicht des Graphen. Kapitel~8 macht aus diesem Unterschied ein ganzes Thema.
\end{warnung}

\begin{warnung}
\boxtitelr{Zweitens: \enquote{Offensichtlich} ist hier besonders gefährlich}
Weil Graphen anschaulich sind, verleiten sie dazu, ein Bild für einen Beweis zu halten. Das geht sehr oft gut und gelegentlich furchtbar schief.

Ein Beispiel, das Sie in Kapitel~7 wiederfinden: \kursiv{Jeder zusammenhängende Graph, in dem jede Ecke mindestens zwei Nachbarn hat, enthält einen Kreis} --- das ist wahr und wird bewiesen. Und: \kursiv{Jeder zusammenhängende Graph, in dem jede Ecke mindestens zwei Nachbarn hat, enthält einen Kreis durch alle Ecken} --- das ist falsch, sieht aber genauso plausibel aus.

\begin{center}
\begin{tikzpicture}[scale=0.9]
\node[knoten] (z) at (0,0) {};
\foreach \a in {90,210,330} {
  \node[knoten] (p\a) at (\a:1.25) {};
  \node[knoten] (q\a) at ($(\a:1.25)+(\a+90:0.6)$) {};
  \node[knoten] (r\a) at ($(\a:1.25)+(\a-90:0.6)$) {};
  \draw[kante] (z) -- (p\a);
  \draw[kante] (p\a) -- (q\a);
  \draw[kante] (p\a) -- (r\a);
  \draw[kante] (q\a) -- (r\a);
}
\node[font=\scriptsize\itshape, color=grau] at (0,-2.35) {jede Ecke hat $\ge 2$ Nachbarn --- aber kein Kreis trifft alle zehn Ecken};
\end{tikzpicture}
\end{center}
\end{warnung}

\begin{warnung}
\boxtitelr{Drittens: Die Bezeichnungen sind uneinheitlich}
Die Graphentheorie ist aus mehreren Quellen zusammengeflossen, und das hört man ihrer Sprache an. Dieselbe Sache heißt je nach Buch anders.

\begin{center}
{\small
\renewcommand{\arraystretch}{1.42}
\begin{tabular}{@{}P{2.6cm}P{3.0cm}P{2.6cm}@{}}
\toprule
\textbf{Hier} & \textbf{Auch üblich} & \textbf{Englisch}\\
\midrule
Ecke & Knoten, Punkt & vertex, node\\
Kante & Verbindung, Linie & edge, line\\
Grad & Valenz, Ordnung & degree, valency\\
Kreis & Zyklus & cycle\\
Weg & Pfad, Kantenzug & path\\
Paarung & Matching & matching\\
\bottomrule
\end{tabular}}
\end{center}

\begin{center}
\begin{minipage}{0.88\textwidth}
\kursiv{Dieser Band sagt durchgehend \enquote{Ecke} und \enquote{Kante}. Wenn Sie in einem anderen Buch \enquote{Knoten} lesen, ist dasselbe gemeint. Besonders heikel ist \enquote{Weg}: Manche Bücher erlauben dabei Wiederholungen von Ecken, andere nicht. Kapitel~2 legt das hier ein für alle Mal fest.}
\end{minipage}
\end{center}
\end{warnung}

\section*{Ein Beispiel gleich zu Beginn}

Damit Sie sehen, wohin die Reise geht, hier eine vollständige mathematische Argumentation --- lesbar, bevor Sie irgendetwas gelernt haben.

\begin{merksatz}
\boxtitel{Auf jeder Feier gibt es zwei Personen mit gleich vielen Bekannten}
\textbf{Die Behauptung.} Auf einer Feier mit $n \ge 2$ Personen gibt es stets zwei Personen, die gleich viele der Anwesenden kennen. (\enquote{Kennen} sei gegenseitig: Wenn $A$ die Person $B$ kennt, dann kennt $B$ auch $A$.)

\schritt{Schritt 1 --- Was kann eine einzelne Person überhaupt an Bekanntenzahlen haben?}
Eine Person kennt mindestens $0$ und höchstens $n-1$ der übrigen Anwesenden (sich selbst zählen wir nicht mit). Es gibt also $n$ mögliche Werte:
\[ 0,\ 1,\ 2,\ \ldots,\ n-1 . \]

\schritt{Schritt 2 --- Der scheinbare Widerspruch}
Wir haben $n$ Personen und $n$ mögliche Werte. Es könnte also sein, dass jede Person einen anderen Wert hat --- die Werte würden dann genau aufgebraucht.

\schritt{Schritt 3 --- Warum das nicht sein kann}
Nehmen wir an, alle $n$ Personen hätten verschiedene Bekanntenzahlen. Dann kommt jeder der $n$ Werte $0, 1, \ldots, n-1$ genau einmal vor. Insbesondere gibt es:
\begin{itemize}[leftmargin=1.4em, itemsep=2pt]
\item eine Person $X$ mit dem Wert $0$ --- sie kennt \kursiv{niemanden} der Anwesenden;
\item eine Person $Y$ mit dem Wert $n-1$ --- sie kennt \kursiv{alle} übrigen $n-1$ Anwesenden.
\end{itemize}

\schritt{Schritt 4 --- Der Widerspruch}
$Y$ kennt alle übrigen Anwesenden. Da $X$ nicht dieselbe Person wie $Y$ ist (sie haben ja verschiedene Werte, und $0 \neq n-1$ wegen $n \ge 2$), gehört $X$ zu diesen übrigen. Also kennt $Y$ die Person $X$.

Weil \enquote{Kennen} gegenseitig ist, kennt dann auch $X$ die Person $Y$. Damit kennt $X$ mindestens eine Person --- im Widerspruch zu $X$'s Wert $0$.

\schritt{Schritt 5 --- Schluss}
Die Annahme aus Schritt~3 führt auf einen Widerspruch. Also war sie falsch: Es können \kursiv{nicht} alle Bekanntenzahlen verschieden sein. Somit gibt es zwei Personen mit gleich vielen Bekannten. $\blacksquare$
\end{merksatz}

\begin{merksatz}
\boxtitel{Was an diesem Beweis bemerkenswert ist}
\begin{center}
\renewcommand{\arraystretch}{1.5}
\begin{tabular}{@{}P{3.0cm}P{6.0cm}@{}}
\toprule
\textbf{Kein Rechnen} & Es kommt keine einzige Rechnung vor --- nur genaues Hinsehen.\\
\textbf{Kein Beispiel} & Der Beweis gilt für \kursiv{jede} Feier, für $n=2$ ebenso wie für $n = 10^{9}$.\\
\textbf{Die Voraussetzung zählt} & Bei $n \ge 2$ ist $0 \neq n-1$. Bei $n=1$ wäre $0 = n-1$, und die Aussage wird sinnlos.\\
\textbf{Es war Graphentheorie} & Personen sind Ecken, Bekanntschaften sind Kanten, Bekanntenzahlen sind Grade.\\
\bottomrule
\end{tabular}
\end{center}

\begin{center}
\begin{minipage}{0.88\textwidth}
\kursiv{Die letzte Zeile ist der Punkt. Sie haben soeben Ihren ersten graphentheoretischen Satz bewiesen, ohne ein einziges Fachwort zu benutzen. In Kapitel~1 lernen Sie die Wörter --- der Gedanke bleibt derselbe.}
\end{minipage}
\end{center}
\end{merksatz}

\section*{Wie dieser Band aufgebaut ist}

\begin{merksatz}
\boxtitel{Sieben Teile, fünfzehn Kapitel}
\begin{center}
{\small
\renewcommand{\arraystretch}{1.42}
\begin{tabular}{@{}P{1.0cm}P{3.4cm}P{4.4cm}@{}}
\toprule
\textbf{Teil} & \textbf{Thema} & \textbf{Die Leitfrage}\\
\midrule
I & Der Begriff (1--3) & Was ist ein Graph, und wie schreibt man ihn auf?\\
II & Bäume (4--5) & Was ist die einfachste zusammenhängende Struktur?\\
III & Durchlaufen (6--7) & Kann man alles genau einmal besuchen?\\
IV & Zeichnen und Färben (8--10) & Wann geht es ohne Kreuzungen, mit wie vielen Farben?\\
V & Paarungen und Flüsse (11--12) & Wie viel passt gleichzeitig hindurch?\\
VI & Grenzen und Zufall (13--14) & Wie viele Kanten erzwingen eine Struktur?\\
VII & Aktueller Stand (15) & Was ist heute offen?\\
\bottomrule
\end{tabular}}
\end{center}
\end{merksatz}

Jedes Kapitel folgt demselben Rhythmus, den Sie aus den anderen Bänden kennen: Motivation, Definition, Begründung, Warnkasten, durchgerechnete Aufgabe mit Probe. Die Kästen sind farbig unterschieden --- blau für Merksätze, rot für Warnungen, grün für Geschichte und Herkunft, violett für den aktuellen Forschungsstand, blaugrau für Definitionen.

Und noch ein Wort zur Herkunft der Beispiele. Wo es sich anbietet, stammen sie aus der Klinik: Dienst\-pläne, Kon\-takt\-nach\-ver\-fol\-gung, Or\-gan\-zu\-tei\-lung, Bet\-ten\-be\-le\-gung. Das ist kein Zierrat. Die Graphentheorie ist in der Medizin an mehr Stellen im Einsatz, als man vermutet, und ein Beispiel, das man kennt, trägt einen Beweis weiter als eines, das man sich erst vorstellen muss.

\vspace{10pt}
\begin{flushright}
\itshape\color{grau}
Nichts ist selbstverständlich.\\
Nichts ist trivial.\\
Jeder Schritt wird erklärt.
\end{flushright}
